% Context from chapters-tex/denoising.tex; for mathematical reference, not compilation.
racle inequality of the form
\eq{
 \mathbb E\norm{f_0-\tilde f^q}^2
 \leq C_0(1+\ln N)\left[\sigma^2+\min_{0\leq M\leq N}\left(\norm{f_0-f_{0,M}^{\mathrm{nlin}}}^2+M\sigma^2\right)\right],
}
where $C_0$ is an absolute constant and $f_{0,M}^{\mathrm{nlin}}$ retains the $M$ largest clean coefficients. In particular, if $\norm{f_0-f_{0,M}^{\mathrm{nlin}}}^2\leq CM^{-2\beta}$ for $1\leq M\leq N$, with $\beta>0$, then for $0<\sigma\leq1$,
\eq{\mathbb E\norm{f_0-\tilde f^q}^2\leq C'(1+\ln N)\sigma^{4\beta/(2\beta+1)},}
where $C'$ depends on $C$ and $\beta$.
\end{thm}

The oracle inequality is a standard Gaussian thresholding bound; see~\cite{donoho-shrinkage}. The rate follows by balancing $CM^{-2\beta}$ and $M\sigma^2$, with integer rounding and truncation at $N$ as in Theorem~\ref{thm-linear-denoising}. The added $\sigma^2$ term accounts for residual noise even at the zero signal.

The universal threshold $T=\sigma\sqrt{2\ln N}$ is conservative: it suppresses nearly all pure-noise coefficients. The smaller thresholds illustrated in Figure~\ref{fig-wavthresh-2d} often give better SNR in these experiments.

                                                     
\subsection{Translation Invariant Thresholding Estimators}

  
\paragraph{Translation invariance.}

Let $f\mapsto\tilde f=\Dd(f)$ be a denoising method, and $f_\tau(x) = f(x-\tau)$ be a translated signal or image for $\tau \in \RR^d$ ($d=1$ or $d=2$).
The denoiser is translation invariant on the shift lattice $\De$ if
\eq{ 
	\foralls \tau \in \De, \quad \Dd(f) = \Dd(f_\tau)_{-\tau}
}
where $\De$ is a lattice of $\RR^d$. A denser shift lattice imposes equivariance at finer spatial resolution.
This corresponds to the fact that $\Dd$ commutes with the translation operator.
\begin{center}
	\image{denoising}{.3}{cycle-spin-principle}
\end{center}
Translation invariance on a sufficiently fine set $\De$ prevents the denoiser from favoring particular feature locations. Without it, the same feature may be reconstructed differently depending on where it appears, producing visible artifacts.

For denoising by thresholding 
\eq{
	\Dd(f) = \sum_m S_T( \dotp{f}{\psi_m} ) \psi_m.
}
a sufficient condition is invariance of the basis $\{\psi_m\}_m$ under shifts in $\De$, up to unit-modulus factors:
\eq{
	\foralls m',\ \foralls\tau\in\De,\ \exists m, \: \exists \la \in \CC,  \quad \quad (\psi_{m'})_\tau = \la \psi_m
}
where $|\la|=1$.

On the periodic domain, the Fourier basis is invariant under every shift $\De=\RR^d$ of $[0,1]^d$. The discrete Fourier basis is invariant under all integer shifts $\De=\{0,\ldots,N_0-1\}^d$, where $N=N_0$ counts samples in one dimension and $N=N_0 \times N_0$ counts pixels in two dimensions.

Unfortunately, an orthogonal wavelet basis 
\eq{
	\{\psi_m = \psi_{j,n} \}_{j,n}
}
is generally not invariant under arbitrary shifts in either the continuous or discrete setting. For instance, in 1-D,
\eq{
	(\psi_{j',n'})_{\tau} \notin \{ \psi_{j,n} \} \qforq \tau=2^{j\prime}/2.
}

  
\paragraph{Cycle spinning.}

For a finite group $\De$ of periodic shifts, average the shifted, denoised, and realigned outputs of $\Dd$ to obtain a translation-invariant estimator:
\eql{\label{eq-cycle-spinning}
	\Dd_{\text{inv}}(f) = \frac{1}{|\De|} \sum_{\tau \in \De} \Dd( f_{\tau} )_{-\tau}.
}
One checks that
\eq{ 
	\foralls \tau \in \De, \quad \Dd_{\text{inv}}(f) = \Dd_{\text{inv}}(f_\tau)_{-\tau}
}
To obtain translation invariance at pixel precision for data with $N$ samples, one should use a set of $|\De|=N$ translation vectors. For wavelets, this requires $O(N^2)$ operations.

Figure \ref{fig-wavthresh-2d-soft-ti} applies cycle spinning to hard thresholding in an orthogonal wavelet basis, using the following shifts:
$\De = \{ 0, 1/N_0, 2/N_0, 3/N_0 \}^2$ on an $N_0\times N_0$ image, corresponding to shifts by $\{ 0, 1, 2, 3 \}^2$ pixels.
This requires 16 forward and inverse wavelet-transform pairs. Averaging this subset of shifts does not in general guarantee exact pixel-shift equivariance. In this example, partial averaging reduces sensitivity to shifts, substantially improves the SNR, and suppresses oscillatory artifacts. These artifacts move as $\tau$ changes, so averaging reduces them.

In Figure \ref{fig-wavthresh-ti}, translation-invariant hard thresholding slightly outperforms translation-invariant soft thresholding, reversing their ranking in the orthogonal-basis experiment.                                                                                               

\myfigure{
\tabdeux{
\image{denoising}{.3}{wavthresh-2d-soft}&
\image{denoising}{.3}{wavthresh-2d-hard-ti}\\
21.8dB & 23.4dB
}
}{ 
	Soft thresholding in an orthogonal wavelet basis (left) and translation-invariant hard thresholding (right).  
}{fig-wavthresh-2d-soft-ti}


\myfigure{
\image{denoising}{.5}{wavthresh-ti-curves}
}{ 
	SNR as a function of $T/\si$ for translation-invariant thresholding.  
}{fig-wavthresh-ti}


                                      
                                                                                                                                                                                                                         

  
\paragraph{Translation invariant wavelet frame.}

Cycle spinning also has a frame interpretation: replace the orthogonal basis $\Bb = \{\psi_m\}$ by its translated copies,
\eql{\label{eq-ti-wavframe}
	\Bb_{\text{inv}} = \{ (\psi_m)_\tau \}_{m, \tau \in \De}.
}
Although $\Bb_{\text{inv}}$ is no longer an orthogonal basis, orthonormality of each translated copy gives the energy and reconstruction identities
\eq{
	\norm{f}^2 = \frac{1}{|\De|} \sum_{m, \tau \in \De} |\dotp{f}{(\psi_m)_\tau}|^2
	\qandq
	f = \frac{1}{|\De|} \sum_{m, \tau \in \De} \dotp{f}{(\psi_m)_\tau} (\psi_m)_\tau.
}
Such redundant families are called tight frames.

One can then define a translation invariant thresholding denoiser
\eql{\label{eq-ti-thresh}
	\Dd_{\text{inv}}(f) = \frac{1}{|\De|} \sum_{m, \tau \in \De} S_T(\dotp{f}{(\psi_m)_\tau}) (\psi_m)_\tau.
}
This estimator equals the cycle spinning estimator defined in \eqref{eq-cycle-spinning}.

Counted with multiplicity, the translated frame has $|\De| |\Bb|$ elements.
For a basis of signals with $N$ samples and a lattice of $|\De|=N$ shifts, this gives up to $N^2$ vectors in $\Bb_{\text{inv}}$. In a hierarchical construction such as a wavelet basis, however, different shifts can produce the same atom:
\eq{
	(\psi_m)_\tau = (\psi_{m'})_{\tau'}
	\qforq
	m \neq m'
	\qandq
	\tau\neq \tau',
}
The number of distinct vectors in $\Bb_{\text{inv}}$ can therefore be much smaller than $N^2$.
For instance, for an orthogonal wavelet basis, one has
\eq{
	(\psi_{j,n})_{k 2^j} = \psi_{j,n+k},
}
so many translated atoms coincide. For a signal of length $